\section{A one-sided coupling estimate}\label{sec:coupling}

The next lemma controls mean square displacement from an upper bound in
only one direction. Large downward displacements are permitted, so
bounding the positive and negative displacements separately would lose
information. Equality of the marginal laws instead balances the expected
change of any integrable function applied to the two variables. We use
an increasing function whose increments dominate squared displacement;
its increments in the bounded direction can be estimated using the common
variance. Neither independence nor exchangeability is assumed.

\Needspace{5\baselineskip}
\begin{lemma}\label{lem:coupling}
Let $U,V$ be real random variables with the same distribution, finite mean
$m$, and finite variance $\sigma^2$. If $V-U\le a$ almost surely, where
$a\ge0$, then
\[
 \E(U-V)^2\le 8a\sigma.
\]
\end{lemma}

\begin{proof}
The function $|t-m|$ grows linearly away from the common mean. Its
primitive is the increasing function
\[
 g(t)=\int_m^t\abs{z-m}\,dz=\tfrac12(t-m)\abs{t-m}.
\]
For $u<v$, we have
\begin{equation}\label{eq:g-increments}
 \frac{(v-u)^2}{4}\le g(v)-g(u)
 \le \frac{v-u}{2}\bigl(\abs{u-m}+\abs{v-m}\bigr).
\end{equation}
For the lower bound, if $u\le m\le v$, the integral is
$\bigl((m-u)^2+(v-m)^2\bigr)/2\ge(v-u)^2/4$; if the interval lies on one
side of $m$, the integral is at least $(v-u)^2/2$.
For the upper bound, convexity places $\abs{z-m}$ below the chord joining
its endpoint values, whose integral is the displayed upper bound.

The second-moment assumption makes $g(U)$ and $g(V)$ integrable.
Their expectations agree. For the integrable random variable
$Y=g(V)-g(U)$, write $Y_+=\max\{Y,0\}$. The identity $\E Y=0$ implies
$\E|Y|=2\E Y_+$. Since $g$ is strictly increasing, $Y>0$ exactly when
$V>U$, and hence
\begin{align*}
 \E\abs{g(U)-g(V)}
 &=2\E\bigl[\mathbf1_{\{V>U\}}(g(V)-g(U))\bigr]\\
 &\le a\E\bigl[\abs{U-m}+\abs{V-m}\bigr]
 \le 2a\sigma.
\end{align*}
The first inequality uses \eqref{eq:g-increments} and $V-U\le a$ on
$\{V>U\}$; the last is Cauchy--Schwarz.
The lower bound in \eqref{eq:g-increments}, with the endpoints put in
increasing order, gives $(U-V)^2\le4|g(U)-g(V)|$ pointwise. Taking
expectations proves the lemma.
\end{proof}


The cancellation $\E[g(V)-g(U)]=0$ uses only equality of marginal laws.
Primitive differences with this property also appear in R\"ollin's
formulation of Stein's method~\cite[Section~1]{Rollin2008}.
The present one-sided second-moment estimate follows from the elementary
increment bounds above.
